\scanpage{115}
% archive sheet number: 114
Considérons le groupe cartographique
\[
 \Gamma=\{\sigma_0,\sigma_1,\sigma_2\mid\sigma_0^2=\sigma_1^2=\sigma_2^2=(\sigma_0\sigma_2)^2=1\}\tag*{(1)}
\]
donc engendré par des gén. $\sigma_0,\sigma_1,\sigma_2$ satisfaisant
\[
 \left\{\begin{aligned}
 \sigma_0^2=\sigma_1^2=\sigma_2^2&=1\\
 \sigma_0\sigma_2&=\sigma_2\sigma_0
 \end{aligned}\right.\tag*{(2)}
\]
\deleted{On pose,} Pour toute partie finie \inserted{non vide} \deleted{\unsure{la plus}} $H$ de $\{0,1,2\}$ \inserted{$H\ne\varnothing\quad H\subset\{0,1,2\}$} on pose
\[
 \left\{\begin{aligned}
 \Gamma_H&=\text{ss-gpe engendré par les }\sigma_i\text{ avec }i\notin H\\
 \Phi_H&=\Gamma/\Gamma_H\qquad(\text{espace homogène})
 \end{aligned}\right.\tag*{(3)}
\]
\deleted{donc}

donc
\[
 H\subset K\ \Rightarrow\ \Gamma_K\subset\Gamma_H\ \Rightarrow\ \Phi_K\xrightarrow{\varphi_{HK}}\Phi_H\tag*{(4)}
\]
On appelle drapeaux universels \uline{de type $H$} les él. de $\Phi_H$, et \deleted{\unsure{l'}} \inserted{$\Phi_H$ sera l'}\unsure{ens.} des \uline{dr. univ. de type $H$}. L'ens.
\[
 \Phi_{\unsure{**|\bullet *}}=\coprod_H\Phi_H\tag*{(5)}
\]
est l'ens des drapeaux, il est ordonné par la relation
\[
 x\in\Phi_H,\ y\in\Phi_K\qquad x\unsure{\preceq|\leq}y\ \Longleftrightarrow\ H\subset K,\ x=\varphi_{HK}y\tag*{(6)}
\]
ce qui est bien une relation d'ordre. Les él. minimaux sont les éléments de
\[
 \Phi_{\{0\}}\overset{\mathrm{dfn}}=\Phi_0,\quad
 \Phi_{\{1\}}\overset{\mathrm{dfn}}=\Phi_1,\quad
 \Phi_{\{2\}}\overset{\mathrm{dfn}}=\Phi_2\tag*{(7)}
\]
appelés resp. sommets, arêtes et faces \uline{universelles}. Sur l'ens des drapeaux minimaux, on a la fonction dim
\[
 \Phi_* =\coprod_{i\in\{0,1,2\}}\Phi_i\xrightarrow{\dim}\{0,2\}\qquad\dim(\Phi_i)\in\{i\}\tag*{(8)}
\]
